\scanpage{102}

Pour (24), on note que

\[\tau_i\tau_{i+1}-\underset{\rho\in k^*}{\rho}\tau_{i+1}\tau_i=\underset{\substack{\rotatebox{-90}{$\in$}\\ k^*}}{\underline{\langle a_{i+1},a'_i\rangle}}\,a'_{i+1}\otimes a_i-\rho\underset{\substack{\rotatebox{-90}{$\in$}\\ k^*}}{\langle a_i,a'_{i+1}\rangle}\,a'_i\otimes a_{i+1}\]

\deleted{qui \unsure{reste|doit} est $\ne0$ car \unsure{d'après les hyp.} la famille de $\{a'_{i+1}\otimes a_i,a'_i\otimes a_{i+1}\}$ est \uline{libre} dans $\operatorname{End}(\unsure{E|\mathcal E})$ --- son image par $\operatorname{End}(E)\longrightarrow E$, $u\longmapsto u(x)$ est formée de $\{a'_{i+1}(x)a_i,a'_i(x)a_{i+1})$, qui \unsure{son} est libre si}

Son image est \unsure{formé|formée} des

\[\langle a_{i+1},a'_i\rangle a'_{i+1}(\unsure{n|x})a_i-\rho\langle a_i,a'_{i+1}\rangle a'_i(x)a_{i+1},\]

je dis que c'est $k(a_i,\unsure{a_{i+1}|h_{i+1}})=h_i+h_{i+1}$, \deleted{ce} qui \deleted{\unsure{provient de|prouve}} signifie aussi que

\[(a'_i,a'_{i+1}):E\unsure{\text{---}|\longrightarrow} k^2=ke_1\oplus ke_2\]

est surjectif, en effet \deleted{l'image de} \inserted{sa restriction} $k(a_{i-1},\deleted{\unsure{a|\oplus}}a_i)$ a la matrice \deleted{$a_{i-1}$ est $(\underset{\substack{\rotatebox{-90}{$\in$}\\ k^*}}{\langle a_{i-1},a'_i\rangle},0)$, donc l'image de $E$ contient $ke_1$, enfin l'image}

\[\begin{pmatrix}\langle a_{i-1},a'_i\rangle&*\\0&\langle a_i,a'_{i+1}\rangle\end{pmatrix}\]

de déterminant $\langle a_{i-1},a'_i\rangle\langle a_i,a'_{i+1}\rangle\in k^*$, en vertu de (25), lui-$=$ \unsure{traduction|conséquence} évidente de (13) $\lambda_n\in k^*$, compte tenu de (4). Enfin (25$'$) est une conséquence immédiate de (25),

cqfd.

% archive sheet number: 101
